\documentclass[12pt,a4paper]{article}
\usepackage{fontspec}
\setmainfont{Liberation Serif}
\usepackage[bulgarian]{babel}
\usepackage[a4paper,margin=2.5cm]{geometry}
\usepackage[expansion=false]{microtype}
\usepackage{amsmath}
\usepackage{graphicx}
\usepackage{booktabs}
\usepackage{array}
\newcolumntype{L}[1]{>{\raggedright\arraybackslash}p{#1}}
\newcolumntype{C}[1]{>{\centering\arraybackslash}p{#1}}
\usepackage{longtable}
\PassOptionsToPackage{hyphens}{url}
\usepackage{hyperref}
\usepackage{xcolor}
\usepackage{enumitem}
\usepackage{listings}
\hypersetup{colorlinks=true,linkcolor=black,citecolor=black,urlcolor=black,filecolor=black}
\setlength{\parindent}{1.25cm}
\setlength{\parskip}{0.25em}
\lstset{basicstyle=\ttfamily\small,breaklines=true,frame=single}

\begin{document}

\begin{titlepage}
\centering
\large Софийски университет „Св. Климент Охридски“\\
Факултет по математика и информатика\\
Катедра „Компютърна информатика“\\[2cm]
\rule{\textwidth}{0.4pt}\\[2cm]
{\LARGE ДИПЛОМНА РАБОТА}\\[1cm]
{\Large Система за извличане с генериране (RAG) за достъп до фармакотерапевтични ръководства на български език}\\[2cm]
Цветелина Пашовска\\
Магистърска програма „Изкуствен интелект“, фак. номер 7MI3400669\\[2cm]
Ръководител:\\
проф. д-р Мария Нишева\\
катедра „Компютърна информатика“, ФМИ, СУ „Св. Климент Охридски“\\
\vfill
София, 2026
\end{titlepage}

\tableofcontents
\newpage

\section{Увод}

\subsection{Достъп до медицинска информация на български език и ролята на извличането с генериране}
Фармакотерапевтичните ръководства, публикувани от Националния съвет по цени и реимбурсиране на лекарствените продукти, съдържат клинично и нормативно значима информация за десетки медицински области, но са структурирани като обемни PDF документи с разнородно вътрешно оформление. За здравен специалист, който търси конкретна препоръка, като дозировка, противопоказание или алгоритъм на лечение, това на практика означава ръчно прелистване на десетки страници, дори когато отговорът вече е формулиран в официално одобрен документ. Проблемът тук не е липса на информация, а нейната труднодостъпност в момента на клиничното решение.

Система от типа извличане с генериране (Retrieval-Augmented Generation, RAG) предлага пряко решение на този проблем. Вместо да разчита единствено на параметричната памет на езиков модел, която може да бъде остаряла или неточна за специфична за България нормативна база, системата извлича релевантни откъси от актуалните ръководства и генерира отговор въз основа на тях, придружен от изричен източник: документ, страница и раздел. Именно възможността отговорът да бъде проследен до конкретен официален документ прави подхода особено подходящ за медицински контекст, в който грешен или непроверим отговор носи реален риск, а не само неудобство.

Настоящата дипломна работа надгражда резултатите от преддипломния проект, в който бе изграден и структуриран корпус от 1380 откъса от 17 национални фармакотерапевтични ръководства, готов за индексиране, като всеки откъс носи данни за произход (provenance): документ, страници и раздел. След преглед на научната литература (arXiv, ACL Anthology, RANLP) и общодостъпни източници на български език не беше открита публикувана система от типа извличане с генериране, насочена конкретно към медицински или фармакотерапевтични ръководства на български език. Намерените български публикации по темата са предимно общи, популярни разяснения на подхода, а не описания на реализирана система върху национален медицински корпус. Тази празнина е допълнителна мотивация за настоящата работа: тя не само прилага установен подход към нов, специализиран корпус, но го прави за език и предметна област, за които все още липсва пряк аналог. 

\subsection{Цели и задачи на дипломната работа}
Основната цел на тази дипломна работа е да се разработи работещ прототип на система за извличане с генериране за достъп до национални фармакотерапевтични ръководства на български език, надграждайки резултатите от преддипломния проект с компоненти за извличане на релевантна информация и генериране на отговори.

За постигането на тази цел са поставени следните задачи:
\begin{enumerate}[noitemsep]
\item аналитичен обзор на литературата по системи за извличане с генериране, методи за извличане (retrieval), приложения в медицинската област и особености при работа с текст на български език;
\item изграждане на векторно представяне (embeddings) на структурираните в преддипломния проект текстови откъси с многоезичен модел, подходящ за българския език, и индексиране във векторна база от данни;
\item реализация на компонент за извличане на релевантна информация (retrieval) по зададена заявка;
\item проектиране и имплементиране на слой за генериране на отговори на база извлечен контекст (prompt engineering);
\item разработване на демонстрационно уеб приложение за задаване на въпроси и преглед на генерираните отговори;
\item качествен анализ на примерни резултати от системата, включващи въпроси от различни медицински области, и дискусия на наблюдаваните ограничения;
\item формулиране на изводи и препоръки за бъдещо развитие.
\end{enumerate}

\section{Литературен преглед}

\subsection{Архитектура и компоненти на системата за извличане с генериране}
Подходът извличане с генериране (Retrieval-Augmented Generation, RAG) е предложен от Lewis и др. \cite{lewis2020rag} като начин да се съчетаят параметричната памет на голям езиков модел с непараметрична памет под формата на външен, обновяем корпус от документи: вместо да разчита единствено на теглата си, моделът получава по заявка извлечени откъси и генерира отговор въз основа на тях. Тази идея решава едновременно два практически проблема: ограничената и статична памет на езиковите модели и невъзможността отговорът да бъде проверен спрямо конкретен източник. Същевременно обаче тя прехвърля тежестта на верността на отговора върху качеството на извличането, а не само върху генеративния модел.

По-новите систематизиращи прегледи очертават еволюция от тази първоначална, линейна схема, при която заявката преминава последователно през извличане, допълване на подсказката и генериране, към т.нар. модулна архитектура, при която всеки компонент -- индексиране, извличане, преработка на извлечения контекст, генериране -- е независимо заменяем и може да се комбинира с допълнителни стъпки, като преформулиране на заявката или повторно класиране на резултатите \cite{wu2024ragnlpsurvey}. Sharma \cite{sharma2025ragsurvey} систематизира съществуващите варианти в няколко ориентации, насочени съответно към подобряване на извличането, на генеративния модул, към хибридни решения и към устойчивост спрямо шумни или противоречиви извлечени откъси, и посочва качеството на извличането като най-честото тясно място в цялостната система. Като по-нататъшно развитие се откроява агентният подход (agentic RAG), при който системата итеративно решава дали и какво да извлече, вместо да следва фиксирана еднократна последователност \cite{singh2025agenticrag}. За целите на настоящата дипломна работа подобна сложност не е оправдана, тъй като въвежда допълнителна латентност и трудности при проследяване на произхода на отговора, без да е ясно, че ще подобри резултата върху затворен и сравнително добре структуриран корпус от национални ръководства. Затова избраната отправна точка тук е модулна, но не агентна архитектура, съобразена със структурата, зададена от преддипломния проект: индексиране на откъсите, самостоятелен компонент за извличане и отделен слой за генериране, всеки от които може да бъде подобряван независимо.

Оценяването на подобни системи е отделен методологически въпрос: освен показателите за самото извличане (виж \ref{sec:retrieval-methods}), генерираният отговор трябва да бъде проверен за вярност спрямо извлечения контекст, тъй като езиковият модел може да „халюцинира“ твърдения, които не се съдържат в подадените откъси. Malin и др. \cite{malin2024faithfulnessmetrics} систематизират съществуващите показатели за вярност (faithfulness) и отбелязват, че оценяването с друг езиков модел като съдия към момента корелира най-добре с човешката преценка, докато Tamber и др. \cite{tamber2025faithfulness} предлагат постоянно обновяван еталон именно за тази цел. Тези резултати мотивират методиката за качествен анализ в тази дипломна работа (виж раздел \ref{sec:qualitative-eval}) да включва не само проверка на извлечения източник, но и отделна преценка доколко генерираният отговор се придържа само към предоставения контекст.

\subsection{Методи за извличане и модели за векторно представяне}
\label{sec:retrieval-methods}
Класическото разграничение в извличането на информация е между лексикално извличане, което съпоставя точни думи или n-грами (например BM25), и семантично извличане, при което заявката и документите се представят като вектори в общо пространство и се сравняват чрез мярка за близост, най-често косинусова близост, независимо от точното съвпадение на думите. Лексикалното извличане е бързо, обяснимо и не изисква обучение, но пропуска синоними и перифрази. Семантичното извличане улавя смисъла, но зависи изцяло от качеството на модела за векторно представяне (embeddings) и е по-трудно за диагностициране, когато върне грешен резултат. За медицински текст, богат на синоними, съкращения и специализирана терминология, например различни търговски и генерични названия на едно и също лекарство, семантичното извличане изглежда по-подходящо. Именно поради тази специализирана терминология обаче качеството на векторното представяне не бива да се приема за дадено, а трябва да бъде проверено емпирично (виж раздел \ref{sec:embedding-choice}).

Що се отнася до самите модели за векторно представяне, Luo и др. \cite{luo2024llmdenseretrieval} показват, че използването на голям езиков модел като енкодер за плътно извличане (dense retrieval) подобрява точността спрямо по-малки специализирани енкодери, но с чувствително по-висока изчислителна цена. Този компромис трябва да се претегли спрямо ограничените изчислителни ресурси на дипломна работа и необходимостта системата да отговаря в разумно време при демонстрация. Извличането по правило се комбинира с векторна база от данни, оптимизирана за приблизително търсене на най-близки съседи. Неотдавнашно емпирично изследване на седем широко използвани системи -- FAISS, Qdrant, Milvus, Weaviate, Chroma, pgvector и LanceDB -- показва съществени разлики в компромиса между точност, латентност и консумация на памет в зависимост от индексиращата структура \cite{rashmiks2026vectordb}. Това налага изборът на векторна база в тази дипломна работа (раздел \ref{sec:embedding-choice}) да бъде направен според мащаба на корпуса от преддипломния проект, тоест 1380 откъса, а не според общи препоръки за много по-големи колекции.

\subsection{Приложения в медицинската област}
Систематизиращият преглед на He и др. \cite{he2025biomedrag} обхваща технологиите, наборите от данни и клиничните приложения на извличането с генериране в биомедицината и подчертава повтарящ се проблем: разминаване между способността на системите да извличат релевантна информация и способността им да я обосноват по проверим начин пред клинициста. Скринингов преглед на Miao и др. \cite{miao2025jmirnursing} върху 67 проучвания в областта на медицината и сестринските грижи установява, че по-голямата част от съществуващите системи са насочени към отговаряне на медицински въпроси и клинична подкрепа при вземане на решения. Приложения, изградени върху затворен корпус от официално одобрени ръководства, какъвто е разглежданият в тази работа, остават сравнително слабо застъпени: повечето изследвани системи разчитат на широки, отворени биомедицински корпуси, например PubMed, а не на регулаторно одобрен, ограничен по обхват източник.

Именно тази разлика е показателна за настоящата дипломна работа. Докато извличането от отворен уеб корпус трябва допълнително да преценява достоверността и авторитетността на всеки източник, затвореният корпус от национални фармакотерапевтични ръководства носи авторитетност „по подразбиране“, а основният риск се измества към точността на самото извличане и към честността на генеративния модул да не добавя информация извън предоставения контекст. Li и Zhu \cite{li2026medicalragindustry} формулират именно такива практически препоръки от индустриален опит за медицински системи за извличане с генериране: стриктно ограничаване на генерирания отговор до предоставения контекст, връщане на изричен източник (документ, страница, раздел) към всяко твърдение и изрично сигнализиране, когато не е намерен релевантен откъс. Тези препоръки директно съответстват на изискването за данни за произход, наследено от преддипломния проект, и са отправна точка при проектирането на слоя за генериране на отговори (раздел \ref{sec:generation}), където е обсъдено и доколко могат да бъдат изпълнени от малък генеративен модел.

\subsection{Особености при работа с по-слабо ресурсни езици}
Повечето цитирани дотук резултати са получени и оценени основно върху английски текст, което поставя въпроса доколко изводите им важат за българския език. Съществуват по същество два алтернативни подхода за извличане на информация, когато езикът на корпуса е по-слабо ресурсен. Първият е кросезиково извличане: заявката или корпусът се превежда към високоресурсен език като английски, след което се използва силен английски модел за извличане и генериране. Ranaldi и др. \cite{ranaldi2025crossrag} показват, че подобен превод преди извличането подобрява резултатите върху задачи, изискващи много фактическо знание, включително за по-слабо ресурсни езици. Този подход обаче въвежда допълнителна стъпка на превод, която може да изкриви медицинска терминология, дозировки и съкращения. Рисковете от подобно изкривяване вече бяха отчетени при OCR обработката в преддипломния проект и биха се повторили при машинен превод. Вторият подход е директна работа на изходния език чрез многоезичен или езиково специфичен модел за векторно представяне, без междинен превод. Той запазва оригиналната терминология, но изисква съответният модел действително да представя добре българския език, а не само големите езици в обучаващата му извадка.

По втория подход има насърчителни, но нееднозначни резултати. От една страна, съществуват специализирани български модели от типа BERT и RoBERTa, обучени основно за класификация, разпознаване на именувани обекти и синтактичен анализ \cite{marinova2023bulgariantransformers, paev2024bulgarianllm}. Това показва, че за българския език вече има зряла инфраструктура за представяне на текст, но тези модели не са били оценявани конкретно за задачата извличане чрез сходство на вектори, което представлява отворен въпрос. От друга страна, изследвания върху други по-слабо ресурсни езици показват, че общите многоезични модели за векторно представяне често подценяват съответния език и се нуждаят от допълнително адаптиране. Omotoso и др. \cite{omotoso2025bgem3} дообучават многоезичния модел BGE-M3 за нигерийски по-слабо ресурсни езици с осезаемо подобрение. Nzeyimana и Rubungo \cite{nzeyimana2025kinyacolbert} изграждат специализиран, лексикално обоснован модел за извличане за езика киняруанда вместо да разчитат на генеричен многоезичен модел. Navasardyan и др. \cite{navasardyan2026embeddings} показват, че дори малък обем шумни, синтетично генерирани данни е достатъчен за забележимо подобряване на векторното представяне на по-слабо ресурсен език. Общото между тези резултати е изводът, че годността на дадено многоезично векторно представяне за конкретен по-слабо ресурсен език не бива да се приема по подразбиране, а трябва да се провери емпирично и, при нужда, да се коригира с целево, дори ограничено по обем, дообучение.

За разлика от векторното представяне, при генеративния компонент положението за българския език е по-благоприятно. Alexandrov и др. \cite{alexandrov2024bggpt} от Института по компютърни науки, изкуствен интелект и технологии (INSAIT) представят BgGPT -- семейство отворени модели, получени чрез продължено предобучение на Gemma~2 върху около 100 милиарда токена, от които 85 милиарда на български, съчетано с техника за сливане на модели (branch-and-merge), ограничаваща катастрофалното забравяне на английски и общи умения при адаптирането към нов език. Моделите се разпространяват в няколко размера -- 2.6, 9 и 27 милиарда параметри -- с отворени тегла и готови квантувани версии за локално изпълнение, а по данни на самите автори по-малкият вариант надминава по представяне на български сравними по размер отворени модели. Наличието на специализиран за българския език генеративен модел, който може да се зарежда изцяло локално, прави обосновано генеративният компонент на настоящата система (раздел \ref{sec:generation}) да се базира на BgGPT, вместо на по-голям, но неспециализиран за българския език модел или на модел, достъпен само през платено API.

От направения преглед следват две методологически решения за настоящата дипломна работа. Първо, вместо предварително да се избере един модел за векторно представяне въз основа само на общите му многоезични показатели, ще бъдат съпоставени емпирично няколко кандидат-модела върху корпуса от преддипломния проект (раздел \ref{sec:embedding-choice}), а директният подход на изходния език ще бъде предпочетен пред превод към английски, за да се запази медицинската терминология непроменена. Второ, за генериране на отговори (раздел \ref{sec:generation}) ще бъде използван BgGPT, зареден локално чрез Ollama, което съответства едновременно на изискването системата да работи без връзка към платено API и на установената по-добра пригодност на този модел за българския език.

\section{Методология}

\subsection{Изходни данни}
Корпусът е резултат от преддипломния проект и се състои от 1380 откъса (chunks), получени чрез разбиване на текста на фрагменти (chunking) от 17 национални фармакотерапевтични ръководства. Всеки откъс запазва данни за произход: документ, диапазон страници, секция и идентификатор (\verb|chunk_id|).

\subsection{Изграждане на векторно представяне и индексиране}
\label{sec:embedding-choice}
В съответствие с извода от литературния преглед (раздел \ref{sec:retrieval-methods}), моделът за векторно представяне не е избран предварително, а се съпоставят три отворени, многоезични кандидата, поддържащи българския език, всеки изпълняван локално чрез библиотеката \nolinkurl{sentence-transformers} без връзка към платено API: \nolinkurl{BAAI/bge-m3} \cite{chen2024bgem3}, \nolinkurl{intfloat/multilingual-e5-base} \cite{wang2024multilinguale5} и \nolinkurl{sentence-transformers/paraphrase-multilingual-mpnet-base-v2} \cite{reimers2020multilingual}. При моделите от семейството e5 текстовете се представят с префиксите \texttt{query:} и \texttt{passage:}, изисквани от авторите им за коректна работа, докато другите два модела не се нуждаят от такъв префикс -- разлика, която се третира изрично в имплементацията (\nolinkurl{src/medguiderag/indexing/embed.py}), за да не се получи подвеждащо ниска прилика само поради пропусната конвенция.

За векторна база е избран FAISS, а не сървър-базирано решение като Qdrant или Milvus, разгледани в раздел \ref{sec:retrieval-methods}: при корпус от 1380 откъса точното (а не приближено) търсене чрез \texttt{IndexFlatIP} върху нормализирани вектори -- еквивалентно на косинусова близост -- отнема пренебрежимо малко време, така че компромисите между точност и латентност при приближено търсене на най-близки съседи, обсъдени по-горе, тук не се проявяват. За всеки кандидат-модел се изгражда отделен индекс с придружаващи метаданни (документ, страници, секция, текст на откъса), което позволява последващо сравнение един срещу друг върху едни и същи заявки (раздел \ref{sec:qualitative-eval}).

Таблица~\ref{tab:embedding-stats} обобщава резултата от индексирането на трите кандидата върху локалната машина (на централния процесор, ограничен до две нишки на процес, за да се избегне изчерпване на паметта). Времето за изчисляване на векторните представяния е еднократен, офлайн разход при построяването на индекса и не се отразява на латентността при отговор на заявка, но е практически показателен за скоростта на итериране по време на разработка: \nolinkurl{BAAI/bge-m3} отнема близо два часа за целия корпус, близо 20 пъти повече от по-леките кандидати. Разликата произтича от по-големия му размер (568 милиона параметри спрямо около 278 милиона за другите два) и по-високата размерност на изхода.

\begin{table}[ht]
\centering
\caption{Резултати от индексирането на трите кандидат-модела върху корпуса от 1380 откъса}
\label{tab:embedding-stats}
\begin{tabular}{L{0.42\textwidth}C{0.13\textwidth}C{0.15\textwidth}C{0.17\textwidth}}
\toprule
Модел & Измерения & Размер на индекса & Време за индексиране\\
\midrule
paraphrase-multilingual-mpnet-base-v2 & 768 & 4.0 MB & 5.9 мин\\
multilingual-e5-base & 768 & 4.0 MB & 15.9 мин\\
BAAI/bge-m3 & 1024 & 5.4 MB & 115.9 мин\\
\bottomrule
\end{tabular}
\end{table}

Разликата във времето за индексиране сама по себе си не определя кой модел е по-подходящ; това ще бъде установено едва чрез качеството на извличането върху ръчно анотирани въпроси (раздел \ref{sec:qualitative-eval}). Тя е основание обаче изборът между кандидатите да отчита и практическата цена на преизчисляване на индекса, ако корпусът или стратегията за фрагментиране се променят по-нататък в разработката. След преминаването към изчисления на графичния процесор във втората итерация (раздел~\ref{sec:second-iteration}) индексирането с \nolinkurl{BAAI/bge-m3} отнема 84 секунди, така че тази цена престава да бъде практическо ограничение.

Трите кандидата се различават и по характеристика, която не личи от таблица~\ref{tab:embedding-stats}, но се оказва решаваща: максималната дължина на входа, която моделът обработва при изчисляване на векторното представяне. Тя е 8192 токена за \nolinkurl{BAAI/bge-m3}, 512 за \nolinkurl{intfloat/multilingual-e5-base} и едва 128 за \nolinkurl{paraphrase-multilingual-mpnet-base-v2}, а по-дългите текстове се отрязват без предупреждение. Откъсите от преддипломния проект са с дължина около 420 думи, т.е. средно около 860 токена, като размерът им е избран преди да е известен моделът за векторно представяне. Последиците от това несъответствие са измерени в раздел~\ref{sec:generation-experiments} и показват, че сравнението на трите кандидата при сегашния размер на откъсите не е равностойно.

\subsection{Компонент за извличане}
\label{sec:retrieval-component}
Компонентът за извличане приема текстова заявка и връща най-много $k$ откъса заедно с оценка на релевантността и пълните данни за произход. В сегашния си вид той е хибриден и се състои от три последователни стъпки (\nolinkurl{src/medguiderag/retrieval/hybrid.py}):
\begin{enumerate}[noitemsep]
\item \emph{Кандидати.} Заявката се търси по два независими начина: семантично, чрез векторното представяне и индекса от раздел~\ref{sec:embedding-choice}, и лексикално, чрез функцията за класиране BM25 \cite{robertson2009bm25} върху същите откъси. Преди лексикалното търсене всяка дума се съкращава до първите си шест знака, за да съвпадат различните граматични форми на една и съща дума. Всеки от двата метода връща по 50 откъса.
\item \emph{Сливане.} Двата списъка се обединяват чрез сливане по реципрочни рангове (Reciprocal Rank Fusion) \cite{cormack2009rrf}: всеки откъс получава сумата от $1/(60+r)$ за всеки от списъците, в които се среща на позиция $r$, и за следващата стъпка се запазват първите 20.
\item \emph{Повторно класиране (reranking).} Двадесетте кандидата се подреждат наново от модела \nolinkurl{BAAI/bge-reranker-v2-m3}, изграден върху \nolinkurl{BAAI/bge-m3}. За разлика от модела за векторно представяне, той обработва заявката и откъса заедно, като кръстосан енкодер (cross-encoder) \cite{nogueira2019passage}, и връща оценка на релевантността между 0 и 1. Запазват се най-много $k$ откъса с оценка поне 0.1. Ако нито един откъс не достига този праг, системата не генерира отговор, а съобщава, че в ръководствата не е открита информация по въпроса.
\end{enumerate}

Всяка от стъпките отговаря на конкретна слабост на по-простите алтернативи, измерена във втората итерация (раздел~\ref{sec:second-iteration}). Лексикалното търсене е добавено, защото двата метода грешат при различни въпроси и обединението им съдържа верния откъс по-често от всеки поотделно. За съкращаването на думите е предпочетено грубо отрязване пред морфологичен анализатор за български, защото не изисква допълнителен езиков ресурс, а вече подобрява съществено лексикалното търсене. Сливането използва само позициите в списъците, а не самите оценки, защото оценките на BM25 и косинусовата близост са в несравними скали. Претеглена сума от двете оценки би изисквала настройване на тегло върху данни, каквито в достатъчен обем липсват, докато константата 60 е стойността, препоръчана от Cormack и др., и не е настройвана. Кръстосаният енкодер е по-точен от сравнението на независимо изчислени вектори, защото вижда заявката и откъса едновременно, но е и много по-бавен: оценяването на всичките 1380 откъса за всяка заявка би било непрактично, затова той се прилага само върху малък брой кандидати. Оценката му е в една и съща скала за всички заявки и затова служи и като праг за отказ. Така изискването за изрично съобщение при липса на релевантна информация (раздел~\ref{sec:generation}) се изпълнява на ниво извличане, а не се поверява на генеративния модел.

Първоначалният вариант -- само семантично търсене и $k$ най-близки откъса (\nolinkurl{src/medguiderag/retrieval/retriever.py}) -- е запазен като алтернативен режим, който се избира от интерфейса. Той служи за отправна точка при сравненията и е единственият практичен режим, когато няма графичен процесор, тъй като на централния процесор повторното класиране отнема около 70 секунди на въпрос. И в двата режима всеки от моделите за векторно представяне от раздел~\ref{sec:embedding-choice} се използва през един и същ интерфейс, което позволява сравняването им без промяна на останалата част от системата.

\subsection{Генериране на отговори}
\label{sec:generation}
За генеративния компонент е избран BgGPT (раздел \ref{sec:retrieval-methods}) във варианта с 2.6 милиарда параметри и 4-битово квантуване, зареден локално чрез Ollama, без връзка към платено API. По-големите варианти (9 и 27 милиарда параметри) изискват повече памет от наличните 6 GB на видеокартата, а 5-битовото квантуване на по-малкия модел не успя да бъде заредено изцяло на нея, въпреки че формално свободната видеопамет е достатъчна.

Подсказката към модела е съставена по следния начин, като всяко от решенията е резултат от експериментите в раздел~\ref{sec:generation-experiments}:
\begin{itemize}
\item Инструкцията и откъсите се подават в едно потребителско съобщение, а инструкцията стои след откъсите, непосредствено преди въпроса. Шаблонът за диалог (chat template) на модела Gemma~2, върху който е изграден BgGPT, няма отделна системна роля, и системното съобщение се подава като допълнителен потребителски ход в самото начало, най-далеч от въпроса.
\item Откъсите се подреждат в обратен ред на релевантността, така че най-близкият до заявката откъс е последен, непосредствено до въпроса.
\item Текстът на откъсите се нормализира: преносите на редове от PDF документите, които при двустранно подравнените редове отделят дори единични думи, се заменят с интервали.
\item Инструкцията изисква отговор само с информация от откъсите, без добавени факти, но не изисква цитиране на източници. Източниците се показват от системата детерминистично, от данните за произход на извлечените откъси, тъй като моделът измисля цитати, когато бъде помолен да ги посочи.
\item Контекстният прозорец (context window) е зададен изрично на 8192 токена, колкото е естественият за Gemma~2, вместо стойността по подразбиране на Ollama от 4096 токена, при която подсказката се отрязва без предупреждение. Поради същото ограничение броят на подаваните откъси е най-много пет.
\item Температурата на генериране, т.е. параметърът, който управлява степента на случайност при избора на следващ токен, е нула. Така един и същ въпрос върху едни и същи откъси получава един и същ отговор, което е предпоставка за сравнимост на експериментите.
\end{itemize}

Спрямо препоръките на Li и Zhu \cite{li2026medicalragindustry} това е съзнателно отклонение. Вместо моделът да посочва източник към всяко твърдение, проследимостта се постига на ниво система: потребителят вижда точните откъси, върху които е изграден отговорът, заедно с документа и страниците им. Изискването за изрично съобщение при липса на релевантна информация също не е поверено на модела, защото при опит да бъде изпълнено чрез инструкция моделът отказва да отговори и когато откъсите съдържат отговора. Решението дали изобщо има релевантен откъс е по-естествено да се взема на ниво извличане и във втората итерация е реализирано чрез праг върху оценката на повторното класиране (раздел~\ref{sec:retrieval-component}).

\section{Експерименти и резултати}

\subsection{Демонстрационно приложение}
% TODO: Задача 5 -- уеб интерфейс за въпроси и преглед на генерираните отговори с
% източници.

\subsection{Първа итерация на подобряване на отговорите}
\label{sec:generation-experiments}
Първите опити с демонстрационното приложение показаха, че отговорите често не възпроизвеждат информацията от показаните източници, а съдържат общи или неподкрепени от тях твърдения. Характерен пример е въпросът „Как се лекува хиперемезис гравидум?“ с модела за векторно представяне \nolinkurl{BAAI/bge-m3}. Извлеченият откъс от Фармакотерапевтичното ръководство по акушерство и гинекология изрежда конкретни групи лекарства: антиеметици (ondansetron, metoclopramide, фенотиазини), антисекреторни средства (esomeprazole, famotidine, ranitidine), инфузии с глюкозни разтвори и електролити, седативни средства. Генерираният отговор споменава инфузионната терапия и антисекреторните средства, но пропуска антиеметиците и добавя твърдения, каквито в откъса няма, например за „намаляване на чернодробната и бъбречната недостатъчност“. Разделът описва експериментите, с които бяха установени причините, и промените, които подобриха отговорите. Промените не са окончателни и разделът завършва с проблемите, които остават за следващата итерация.

\paragraph{Протокол.} За да се отделят отделните причини, експериментите следват един и същ протокол. Използват се три въпроса от различни области, за които корпусът съдържа отговор: хиперемезис гравидум (акушерство и гинекология), кога се прилага инсулин при захарен диабет тип 2 (ендокринология) и лечение на атопичен дерматит (дерматология). За всеки въпрос откъсите се извличат веднъж с \nolinkurl{BAAI/bge-m3} ($k=5$) и остават фиксирани, а между вариантите се променя по един фактор на подсказката. Генерирането е с температура нула, за да се изключи случайността. Всеки отговор се оценява по два начина. Автоматично се брои колко от предварително изброените факти в релевантния откъс се срещат в отговора (съответно 14, 7 и 8 факта, най-често имена на лекарства и мерки). Ръчно отговорът се прочита изцяло и се отбелязват твърденията, които не се съдържат в откъсите. Автоматичната мярка е груба, защото отчита само наличието на ключова дума, и при разминаване предимство има ръчният прочит.

\paragraph{Отрязване на подсказката.} Първата находка не е свързана с формулировката на подсказката, а с размера ѝ. Пет откъса заедно с инструкцията заемат между 5500 и 6300 токена, а Ollama по подразбиране използва контекстен прозорец от 4096 токена. При препълване библиотеката llama.cpp, чрез която Ollama изпълнява модела, запазва само последната половина от прозореца. Така подсказката е била съкратена до около 2000 токена, без предупреждение, като отпадат първите откъси, които при подреждане по релевантност са и най-важните. Отрязването е установено по броя обработени токени, който Ollama отчита за всяка заявка. Контекстният прозорец е увеличен до 8192 токена, а броят на откъсите е ограничен до пет.

\paragraph{Системна роля, цитиране и отказ.} Първоначалната подсказка съдържаше правилата в системно съобщение: да се използва само предоставеният контекст, да се цитира източник във формат „[документ, стр.]“ и изрично да се съобщава, когато липсва информация. Тъй като шаблонът за диалог на Gemma~2 няма системна роля, правилата попадат в отделен потребителски ход в самото начало, а при описаното по-горе отрязване на подсказката те са отпадали изцяло. Затова правилата бяха преместени непосредствено преди въпроса, а контекстният прозорец беше увеличен. Проверката на този вариант с откъси, извлечени с \nolinkurl{paraphrase-multilingual-mpnet-base-v2}, разкри два по-сериозни проблема. С дословно зададена фраза за отказ моделът отговаря с нея на всички въпроси, включително на въпроса за инсулина, чийто отговор се съдържа във втория извлечен откъс. Без тази фраза моделът отговаря, но измисля източници: за контролния въпрос „Колко е висок връх Мусала?“, на който корпусът няма отговор, посочва хематологичното ръководство, а в друг вариант несъществуващ файл „ФТР\_География.pdf“. В медицински контекст измисленият източник е по-опасен от липсващия, защото придава привидна достоверност на отговора. Затова цитирането и отказът са изнесени извън генеративния модел (раздел~\ref{sec:generation}).

\paragraph{Варианти на подсказката.} Таблица~\ref{tab:prompt-variants} обобщава вариантите, сравнени при фиксирани откъси. Изходният вариант V0 съответства на състоянието на системата в момента на наблюдавания проблем: инструкцията е след откъсите, а контекстният прозорец е 8192 токена.

\begin{table}[ht]
\centering
\small
\caption{Възпроизведени факти от релевантния откъс при различни варианти на подсказката (откъси, извлечени с \texttt{BAAI/bge-m3}, температура нула)}
\label{tab:prompt-variants}
\setlength{\tabcolsep}{4pt}
\begin{tabular}{C{0.09\textwidth}L{0.39\textwidth}C{0.14\textwidth}C{0.1\textwidth}C{0.11\textwidth}}
\toprule
Вариант & Описание & Хиперемезис & Инсулин & Дерматит\\
\midrule
V0 & изходен: необработен текст, 5 откъса по реда на релевантност & 1/14 & 1/7 & 3/8\\
V1 & V0 с нормализиран текст & 1/14 & 1/7 & 3/8\\
V2 & V1, най-релевантният откъс последен & 9/14 & 2/7 & 4/8\\
V3 & V2 с 3 откъса & 7/14 & 3/7 & 4/8\\
V4 & V2 с конкретизираща инструкция & 7/14 & 5/7 & 4/8\\
V5 & V4 с 3 откъса & 7/14 & 6/7 & 4/8\\
V6 & само най-релевантният откъс, конкретизираща инструкция & 7/14 & 4/7 & 4/8\\
\bottomrule
\end{tabular}
\end{table}

\paragraph{Ред на откъсите.} Най-голямото подобрение идва от реда на откъсите (V1 спрямо V2): при въпроса за хиперемезис възпроизведените факти нарастват от 1 на 9 от 14. При подреждане по релевантност най-подходящият откъс е първи, на около 5000 токена от въпроса, и моделът на практика го пренебрегва, като съставя отговора от последните, несвързани откъси и от собствената си параметрична памет. Твърдението за „намаляване на чернодробната и бъбречната недостатъчност“ не се съдържа в тази форма в нито един от откъсите. То по-скоро съчетава „подобряване на чернодробната функция“ от релевантния откъс с темите на два от несвързаните, които са от ръководствата по гастроентерология и нефрология. Резултатът е в съгласие с наблюденията на Liu и др. \cite{liu2024lostmiddle}, според които езиковите модели използват най-добре информацията в началото и в края на дълъг контекст и най-слабо тази в средата. При BgGPT с 2.6 милиарда параметри ефектът е дори по-силен: в наблюдаваните случаи позицията в началото не е достатъчна, ако между нея и въпроса има хиляди токени несвързан текст, и решаваща е близостта до въпроса.

\paragraph{Нормализиране на текста.} Замяната на преносите на редове с интервали (V0 спрямо V1) не променя броя възпроизведени факти, но съкращава подсказката с около 6\% и премахва накъсването на изреченията, затова е запазена.

\paragraph{Конкретизираща инструкция.} Инструкция, която изрично иска да се изброят мерките, лекарствата и дозировките, посочени в откъсите (V4--V6), повишава автоматичната оценка при въпроса за инсулина, но ръчният прочит показва обратното. Моделът попълва исканата структура с измислено съдържание: при хиперемезис добавя диетичен режим („избягване на мазни и пикантни храни“), какъвто в откъса няма, приписва на лекарствата погрешни ATC кодове и причислява инхибиторите на протонната помпа към антиеметиците, а при инсулина посочва взаимно противоречиви прагове на гликирания хемоглобин. Затова е запазена по-простата инструкция. Случаят илюстрира и ограничението на автоматичната мярка: тя отчита наличието на ключова дума, но не и верността на твърдението около нея.

\paragraph{Брой откъси.} С три откъса (V3) отговорите остават верни спрямо източниците, а подсказката е с около 40\% по-кратка, но отговорът за хиперемезис е по-беден (7 срещу 9 от 14 факта). Запазен е вариантът с пет откъса.

\paragraph{Температура.} При температура 0.2 три последователни изпълнения на въпроса за инсулина върху едни и същи пет откъса (извлечени с \nolinkurl{paraphrase-multilingual-mpnet-base-v2}) дадоха три различни отговора, единият с невярно твърдение: хипогликемията е представена като показание за инсулиново лечение, докато тя е риск от него. Температура нула прави отговорите повторими и е условие за сравнимост на експериментите.

\paragraph{Резултат.} След промените (вариант V2) отговорът за хиперемезис изрежда антиеметиците по групи (ondansetron, metoclopramide, chlorpromazine, promethazine, trifluoperazine), антисекреторните средства (esomeprazole, famotidine, ranitidine, атропин), инфузиите с глюкозни и левулозни разтвори, корекцията на електролитите, поливитамините и седативните средства, като всички тези сведения се съдържат в откъса. Пропуснати остават хоспитализацията, покоят и диетичните препоръки от първоначалните грижи.

\paragraph{Зависимост от извличането.} Промените в подсказката подобряват отговора само когато релевантният откъс е сред извлечените. Таблица~\ref{tab:models-questions} показва, че при въпроса за хиперемезис това е изпълнено единствено за \nolinkurl{BAAI/bge-m3}. Другите два модела не извличат откъса с отговора сред първите пет, а връщат откъси от ръководствата по нуклеарна медицина, нефрология, ендокринология и неврология, и това се запазва и при правилното изписване „гравидарум“ в заявката. При другите два въпроса трите модела извличат релевантни откъси и отговорите са сравними.

\begin{table}[ht]
\centering
\small
\caption{Позиция на откъса с отговора сред петте извлечени и възпроизведени факти по модели за векторно представяне (подсказка V2)}
\label{tab:models-questions}
\setlength{\tabcolsep}{4pt}
\begin{tabular}{L{0.38\textwidth}C{0.2\textwidth}C{0.12\textwidth}C{0.14\textwidth}}
\toprule
Модел & Хиперемезис & Инсулин & Дерматит\\
\midrule
bge-m3 & №1; 9/14 & 2/7 & №1; 4/8\\
multilingual-e5-base & не е извлечен; 2/14 & 3/7 & №1; 3/8\\
paraphrase-multilingual-mpnet-base-v2 & не е извлечен; 0/14 & 2/7 & №2; 4/8\\
\bottomrule
\end{tabular}
\end{table}

\paragraph{Ограничение на дължината на входа.} Причината е несъответствието, посочено в раздел~\ref{sec:embedding-choice}, между дължината на откъсите и максималната дължина на входа на моделите (таблица~\ref{tab:truncation}). Откъсите са с медиана около 860 токена и максимум около 1410. \nolinkurl{paraphrase-multilingual-mpnet-base-v2} представя средно едва 15\% от всеки откъс, т.е. първите около 60 думи, а в откъса за хиперемезис терминът се появява едва около 440-ия токен, така че моделът никога не го „вижда“. \nolinkurl{multilingual-e5-base} обхваща термина, но не и описанието на лечението след него. Освен това стойностите на сходство при него са сгъстени (0.84--0.85 за всичките пет извлечени откъса), така че и малки отклонения в представянето променят класирането. Само \nolinkurl{BAAI/bge-m3} представя откъсите изцяло.

\begin{table}[ht]
\centering
\small
\caption{Обхват на представянето на откъсите при максималната дължина на входа на всеки модел}
\label{tab:truncation}
\setlength{\tabcolsep}{4pt}
\begin{tabular}{L{0.32\textwidth}C{0.12\textwidth}C{0.12\textwidth}C{0.13\textwidth}C{0.15\textwidth}}
\toprule
Модел & Макс. вход (токени) & Отрязани откъси & Представена част & Позиция на „hyperemesis“\\
\midrule
bge-m3 & 8192 & 0\% & 100\% & 438\\
multilingual-e5-base & 512 & 99\% & 60\% & 439\\
paraphrase-multilingual-mpnet-base-v2 & 128 & 100\% & 15\% & 437\\
\bottomrule
\end{tabular}
\end{table}

Следователно досегашното сравнение на трите модела не е равностойно, защото два от кандидатите се оценяват върху непълно представяне на откъсите. Размерът на откъсите е определен в преддипломния проект преди избора на модел за векторно представяне и за \nolinkurl{paraphrase-multilingual-mpnet-base-v2} е неподходящ по конструкция.

\paragraph{Отворени проблеми за следващата итерация.}
\begin{enumerate}[noitemsep]
\item \emph{Несвързани откъси.} Когато по дадена тема корпусът съдържа един релевантен откъс, фиксираният брой $k=5$ добавя четири несвързани. При хиперемезис те са от ръководствата по ендокринология, гастроентерология и нефрология, със сходство около 0.58 срещу 0.63 за релевантния, така че праг, който да ги отдели, не е очевиден. След обръщането на реда те вече не влияят на отговора, но продължават да се показват като източници. Възможни решения са повторно класиране с отделен модел и праг на сходството, калибриран поотделно за всеки модел за векторно представяне, тъй като стойностите на сходство при трите модела са в различни диапазони.
\item \emph{Размер на откъсите спрямо ограничението на моделите.} За равностойно сравнение откъсите трябва да бъдат съобразени с максималната дължина на входа, например чрез по-кратки подоткъси, чиито представяния се съпоставят с изходния откъс.
\item \emph{Граници на откъсите.} Откъсът за хиперемезис завършва по средата на списъка с използвани лекарства, така че част от информацията остава в съседния откъс, който не е извлечен.
\item \emph{Формална оценка.} Изводите в този раздел почиват на три въпроса и следва да бъдат проверени върху по-голям набор от ръчно анотирани въпроси (раздел~\ref{sec:qualitative-eval}).
\item \emph{Ресурси.} Сървърът, чрез който Ollama изпълнява модела, пази в оперативната памет кеш на обработените подсказки с лимит 8 GB. При серия от дълги подсказки кешът достигна този лимит и изчерпа паметта на машината за разработка, така че зареждането на други модели стана невъзможно. Това е практическо ограничение за локалното изпълнение, което следва да се отчита при по-нататъшните експерименти.
\end{enumerate}

\subsection{Втора итерация на подобряване на системата}
\label{sec:second-iteration}
Втората итерация е насочена към отворените проблеми от първата, свързани с извличането: несвързаните откъси, които фиксираният брой $k$ добавя към контекста, неравностойното сравнение на моделите за векторно представяне и липсата на количествена оценка върху повече от три въпроса. За разлика от първата итерация, в която се променяше подсказката при фиксирани откъси, тук се сравняват методи за извличане. Затова първо е съставен набор за разработка с автоматично изчислими показатели. Промяната в извличането разкри и нов проблем при генерирането, описан в края на раздела.

\paragraph{Набор за разработка.} От всяко от 17-те ръководства са избрани случайно по един откъс, а от четирите ръководства с поне 100 откъса -- по два, като изборът е ограничен до откъси, в които се споменават дози. Към тях са добавени трите откъса от първата итерация, общо 24. За всеки откъс е формулиран въпрос, на който той отговаря, и този откъс се счита за верния резултат от извличането. Наборът служи за сравнение на варианти по време на разработката, а не за окончателна оценка, и има две ограничения, които трябва да се имат предвид при тълкуването на резултатите. Първо, при 24 въпроса един въпрос съответства на 0.04 от всеки показател, така че разлики от един-два въпроса не са показателни. Второ, въпросите са формулирани при наличен текст на откъса и неизбежно повтарят част от думите му, което облагодетелства лексикалното търсене. Окончателната оценка върху независимо формулирани въпроси е предмет на раздел~\ref{sec:qualitative-eval}.

\paragraph{Показатели.} Използват се два стандартни показателя за извличане. Пълнотата при $k$ (Recall@$k$) е делът на въпросите, при които верният откъс е сред първите $k$ извлечени. Пълнотата при 5 съответства на броя откъси, подавани на генеративния модел, а пълнотата при 20 -- на броя кандидати за повторното класиране. Средният реципрочен ранг (Mean Reciprocal Rank) е средната стойност на $1/r$, където $r$ е позицията на верния откъс, като се отчита нула, ако той не е сред първите десет. За разлика от пълнотата, този показател отчита и колко високо е класиран верният откъс. Оценката е строга: за верен се приема само конкретният откъс, въпреки че съседните откъси от същото ръководство понякога съдържат част от отговора.

\begin{table}[ht]
\centering
\small
\caption{Сравнение на методите за извличане върху набора за разработка (24 въпроса)}
\label{tab:retrieval-methods}
\setlength{\tabcolsep}{4pt}
\begin{tabular}{L{0.44\textwidth}C{0.1\textwidth}C{0.1\textwidth}C{0.1\textwidth}C{0.13\textwidth}}
\toprule
Метод & Пълнота при 1 & Пълнота при 5 & Пълнота при 20 & Среден реципрочен ранг\\
\midrule
\multicolumn{5}{l}{\emph{Семантично търсене}}\\
paraphrase-multilingual-mpnet-base-v2 & 0.04 & 0.25 & 0.46 & 0.15\\
multilingual-e5-base & 0.42 & 0.62 & 0.75 & 0.50\\
embeddinggemma & 0.33 & 0.71 & 0.83 & 0.48\\
qwen3-embedding (0.6 млрд. параметри) & 0.46 & 0.83 & 0.96 & 0.61\\
bge-m3 & 0.54 & 0.79 & 0.92 & 0.65\\
\midrule
\multicolumn{5}{l}{\emph{Лексикално търсене}}\\
BM25 & 0.25 & 0.79 & 0.92 & 0.44\\
BM25 със съкращаване на думите & 0.42 & 0.79 & 0.92 & 0.55\\
\midrule
\multicolumn{5}{l}{\emph{Хибридно търсене}}\\
сливане на bge-m3 и BM25 & 0.46 & 0.83 & 0.96 & 0.60\\
сливане и повторно класиране & \textbf{0.58} & \textbf{0.92} & \textbf{0.96} & \textbf{0.74}\\
\bottomrule
\end{tabular}
\end{table}

\paragraph{Модели за векторно представяне.} Таблица~\ref{tab:retrieval-methods} потвърждава количествено наблюденията от първата итерация. \nolinkurl{paraphrase-multilingual-mpnet-base-v2} е практически неизползваем: верният откъс е на първо място само при един от 24 въпроса и не е сред първите 20 при повече от половината. Следователно моделът не може да служи дори като първа стъпка, чиито грешки да се поправят от повторното класиране. \nolinkurl{multilingual-e5-base} е по-добър, но отстъпва на \nolinkurl{BAAI/bge-m3} по всички показатели, в съответствие с по-непълното представяне на откъсите (таблица~\ref{tab:truncation}). Като алтернативи на \nolinkurl{BAAI/bge-m3} са изпробвани и два по-нови компактни многоезични модела, изпълнявани чрез Ollama, чиято максимална дължина на входа е достатъчна за откъсите: \nolinkurl{qwen3-embedding} с 0.6 милиарда параметри \cite{zhang2025qwen3embedding} и \nolinkurl{embeddinggemma} с 0.3 милиарда параметри \cite{vera2025embeddinggemma}. За да е сравнението равностойно, \nolinkurl{BAAI/bge-m3} е изпълнен и по същия начин и дава практически същите резултати като през библиотеката \nolinkurl{sentence-transformers}. \nolinkurl{qwen3-embedding} е наравно с \nolinkurl{BAAI/bge-m3}: по-добър е по пълнота при 5 и при 20, по-слаб по пълнота при 1, като разликите са от един-два въпроса. \nolinkurl{embeddinggemma} е по-слаб. Нито един от двата не дава основание за смяна на модела. Семейството Qwen3 включва и модели за повторно класиране, които не са изпробвани в тази итерация.

\paragraph{Варианти на представянето.} Изпробвани са и три начина за подобряване на представянето при \nolinkurl{BAAI/bge-m3}: нормализиране на текста, добавяне на заглавието на ръководството и на раздела пред текста на откъса и разделяне на откъса на застъпващи се подоткъси от 120 или 200 думи, при което откъсът се класира по най-близкия си подоткъс. Последният вариант беше предложен в първата итерация като начин да се съобразят откъсите с ограничението на дължината на входа. Нито един от вариантите не подобрява последователно резултата: пълнотата при 1 е между 0.38 и 0.54, пълнотата при 5 -- между 0.67 и 0.83, а средният реципрочен ранг -- между 0.57 и 0.66, спрямо 0.54, 0.79 и 0.65 за изходния вариант. Разликите са в границите на един-два въпроса и в двете посоки. Вероятното обяснение е, че \nolinkurl{BAAI/bge-m3} и без това обработва откъсите изцяло, а подоткъсите разделят сведения, които в заявката са свързани, например името на заболяването и описанието на лечението. Поради това разделянето е отхвърлено. За двата модела с кратък вход то не е изпробвано, защото при наличието на \nolinkurl{BAAI/bge-m3} подобряването им има само сравнителна стойност.

\paragraph{Лексикално търсене.} BM25 без обработка на думите класира верния откъс на първо място едва при 6 от 24 въпроса. Причината е богатата морфология на българския език: „лечение“, „лечението“ и „лечения“ са различни думи за функцията за класиране. Съкращаването на всяка дума до първите шест знака, грубо приближение на основата ѝ (stemming), повишава пълнотата при 1 от 0.25 на 0.42 и средния реципрочен ранг от 0.44 на 0.55. По пълнота при 5 и при 20 лексикалното търсене се изравнява с \nolinkurl{BAAI/bge-m3}, въпреки че не използва никакъв обучен модел. Това отчасти се дължи на отбелязаното облагодетелстване от начина на съставяне на въпросите. По-важно е обаче, че двата метода грешат при различни въпроси: сред първите 20 резултата \nolinkurl{BAAI/bge-m3} намира верния откъс при два въпроса, при които лексикалното търсене не го намира, и обратно, а поне единият от тях го намира при всичките 24 въпроса. Дължината от шест знака не е настройвана систематично, а морфологичен анализатор би бил по-прецизна алтернатива.

\paragraph{Сливане и повторно класиране.} Сливането по реципрочни рангове повишава пълнотата при 20 до 0.96, т.е. обединеният списък съдържа верния откъс при 23 от 24 въпроса. Сливането обаче понижава пълнотата при 1 спрямо \nolinkurl{BAAI/bge-m3} самостоятелно (0.46 срещу 0.54). То увеличава шанса верният откъс да е сред кандидатите, но не и да е на първо място, тъй като позицията му се усреднява с тази от по-слабия за дадения въпрос метод. Затова сливането е само първа стъпка, последвана от повторно класиране на първите 20 кандидата. Така се постигат най-добрите стойности по всички показатели: пълнота при 1 от 0.58, пълнота при 5 от 0.92 и среден реципрочен ранг от 0.74. Спрямо \nolinkurl{BAAI/bge-m3} самостоятелно най-голямо е подобрението при пълнотата при 5, която определя дали верният откъс изобщо достига до генеративния модел: от 19 на 22 въпроса. Горната граница на подхода е пълнотата на кандидатите: при единствения въпрос, чийто откъс не е сред тях, повторното класиране не може да помогне. Повторното класиране има и ограничение: при въпроса за проследяването на бъбречната функция при хепатит B първите пет откъса са от други ръководства, с оценки от 0.78 до 0.98. Високата оценка показва, че откъсът е тематично близък до въпроса, но не гарантира, че е от вярното ръководство.

\paragraph{Праг за отказ.} Оценката на модела за повторно класиране е в една и съща скала за всички заявки, за разлика от косинусовата близост, чиито стойности зависят от модела за векторно представяне и при \nolinkurl{multilingual-e5-base} са сгъстени в тесен интервал. Затова прагът е определен върху нея. За калибрирането освен 24-те въпроса от набора за разработка са използвани десет въпроса извън обхвата на корпуса. Седем от тях са общи („Колко е висок връх Мусала?“, „Как се приготвя баница?“), а три са близки до медицината, но не и до фармакотерапевтичните ръководства („Как се лекува счупен крак?“, „Как се поставя пломба на зъб?“, „Каква е дозата на парацетамол за куче?“). Най-високата оценка при въпросите извън обхвата е 0.049, при въпроса за парацетамола, а при останалите е под 0.025. Най-ниската оценка на първия откъс при въпросите от набора е 0.717. Прагът 0.1 е избран в интервала между двете стойности, с двукратен запас над най-високата оценка извън обхвата. Така системата отказва да отговори на всичките десет въпроса извън обхвата, без да вика генеративния модел. Прагът има и втори ефект, който решава първия отворен проблем от първата итерация. Прилага се към всеки откъс поотделно, затова отпадат несвързаните откъси, които фиксираният брой $k$ добавяше. При въпроса за хиперемезис остава само релевантният откъс, с оценка 0.98, докато останалите кандидати получават най-много 0.02, при косинусова близост 0.58 срещу 0.63. От 24-те въпроса 21 запазват пет откъса, а три -- по-малко (един, два и три), като верният откъс не е отпаднал при нито един въпрос. Изводът за прага почива на малък брой въпроси извън обхвата, повечето от които са далеч от темата на корпуса. Най-рискови са медицински въпроси по теми, които ръководствата не покриват, но с близък речник, и прагът трябва да бъде проверен върху такива въпроси при окончателната оценка.

\paragraph{Изчислителна цена.} На централния процесор повторното класиране на 20 кандидата отнема около 70 секунди на въпрос, което е неприемливо за интерактивна система. Машината за разработка разполага с графичен процесор NVIDIA GeForce RTX 2060 с 6 GB памет, който дотогава се използваше само от Ollama за генеративния модел, докато библиотеката PyTorch, чрез която се изпълняват моделите за векторно представяне и повторно класиране, беше във вариант само за централен процесор. След преминаването към варианта ѝ за графичен процесор (CUDA) моделите се изпълняват в половинна точност (fp16), т.е. с 16-битови числа с плаваща запетая вместо 32-битови. Това намалява двойно заеманата памет и е необходимо, за да се поберат всички модели в 6 GB. Таблица~\ref{tab:compute} обобщава резултата. Половинната точност не променя резултатите: подредбата след повторното класиране съвпада изцяло с получената на централния процесор, а векторните представяния на откъсите имат косинусова близост поне 0.9997 с изчислените преди това. Затова вече построените индекси са запазени. При заредени три модела -- за векторно представяне, за повторно класиране и генеративен -- заетата видеопамет е около 5.6 от 6 GB, заедно с тази на операционната система, като генеративният модел остава изцяло на графичния процесор.

\begin{table}[ht]
\centering
\small
\caption{Време за изпълнение на централния и на графичния процесор}
\label{tab:compute}
\setlength{\tabcolsep}{4pt}
\begin{tabular}{L{0.5\textwidth}C{0.18\textwidth}C{0.22\textwidth}}
\toprule
Задача & Централен процесор & Графичен процесор, половинна точност\\
\midrule
Повторно класиране на 20 кандидата за един въпрос & около 70 с & 0.97 с\\
Индексиране на корпуса с \texttt{bge-m3} & 116 мин & 84 с\\
Цялото извличане за един въпрос & около 70 с & около 1 с\\
\bottomrule
\end{tabular}
\end{table}

% DECODING-PLACEHOLDER

\subsection{Качествен анализ и оценка}
\label{sec:qualitative-eval}
% TODO: Задача 6 -- примерни въпроси от различни медицински области, Recall@k, MRR,
% точност на данните за произход (provenance accuracy); аналитична дискусия на
% наблюдаваните ограничения.

\section{Техническа реализация}

\subsection{Използвани инструменти и технологии}
% TODO: език, библиотеки за векторно представяне/векторна база/генериране,
% инструментариум (framework) за демонстрационното приложение; кратка обосновка на
% всеки избор спрямо разгледани алтернативи.

\subsection{Структура на проекта и файлове}
% TODO: описание на модулите в \texttt{src/medguiderag/} (\texttt{ingest},
% \texttt{indexing}, \texttt{retrieval}, \texttt{generation}, \texttt{app}) и връзката
% им с преддипломния корпус.

\section{Изводи и препоръки за бъдещо развитие}
% TODO: Задача 7.

\section{Заключение}
% TODO

\begin{thebibliography}{99}

\bibitem{lewis2020rag}
P. Lewis, E. Perez, A. Piktus, F. Petroni, V. Karpukhin, N. Goyal, H. Küttler, M. Lewis, W. Yih, T. Rocktäschel, S. Riedel, D. Kiela.
\emph{Retrieval-Augmented Generation for Knowledge-Intensive NLP Tasks}.
Advances in Neural Information Processing Systems (NeurIPS), 2020. arXiv:2005.11401.

\bibitem{wu2024ragnlpsurvey}
S. Wu, Y. Xiong, Y. Cui, H. Wu, C. Chen, Y. Yuan, L. Huang, X. Liu, T.-W. Kuo, N. Guan, C. J. Xue.
\emph{Retrieval-Augmented Generation for Natural Language Processing: A Survey}.
arXiv:2407.13193, 2024.

\bibitem{sharma2025ragsurvey}
C. Sharma.
\emph{Retrieval-Augmented Generation: A Comprehensive Survey of Architectures, Enhancements, and Robustness Frontiers}.
arXiv:2506.00054, 2025.

\bibitem{singh2025agenticrag}
A. Singh, A. Ehtesham, S. Kumar, T. Talaei Khoei, A. V. Vasilakos.
\emph{Agentic Retrieval-Augmented Generation: A Survey on Agentic RAG}.
arXiv:2501.09136, 2025.

\bibitem{malin2024faithfulnessmetrics}
B. Malin, T. Kalganova, N. Boulgouris.
\emph{A Review of Faithfulness Metrics for Hallucination Assessment in Large Language Models}.
arXiv:2501.00269, 2024.

\bibitem{tamber2025faithfulness}
M. S. Tamber, F. S. Bao, C. Xu, G. Luo, S. Kazi, M. Bae, M. Li, O. Mendelevitch, R. Qu, J. Lin.
\emph{Benchmarking LLM Faithfulness in RAG with Evolving Leaderboards}.
arXiv:2505.04847, 2025.

\bibitem{luo2024llmdenseretrieval}
K. Luo, M. Qin, Z. Liu, S. Xiao, J. Zhao, K. Liu.
\emph{Large Language Models as Foundations for Next-Gen Dense Retrieval: A Comprehensive Empirical Assessment}.
arXiv:2408.12194, 2024.

\bibitem{rashmiks2026vectordb}
A. Rashmiks, T. Madushanka.
\emph{A Comprehensive Empirical Evaluation of Vector Database Systems for Approximate Nearest Neighbor Search: Performance, Quality, and Resource Trade-offs}.
arXiv:2608.12812, 2026.

\bibitem{he2025biomedrag}
J. He, B. Zhang, H. Rouhizadeh, Y. Chen, R. Yang, J. Lu, X. Chen, N. Liu, D. Teodoro.
\emph{Retrieval-Augmented Generation in Biomedicine: A Survey of Technologies, Datasets, and Clinical Applications}.
arXiv:2505.01146, 2025.

\bibitem{miao2025jmirnursing}
Y. Miao, Y. Zhao, Y. Luo, H. Wang, Y. Wu.
\emph{Improving Large Language Model Applications in the Medical and Nursing Domains With Retrieval-Augmented Generation: Scoping Review}.
Journal of Medical Internet Research, 27:e80557, 2025.

\bibitem{li2026medicalragindustry}
L. Li, W. Zhu.
\emph{Pursuing Best Industrial Practices for Retrieval-Augmented Generation in the Medical Domain}.
arXiv:2602.03368, 2026.

\bibitem{ranaldi2025crossrag}
L. Ranaldi, B. Haddow, A. Birch.
\emph{Multilingual Retrieval-Augmented Generation for Knowledge-Intensive Task}.
arXiv:2504.03616, 2025.

\bibitem{marinova2023bulgariantransformers}
I. Marinova, K. Simov, P. Osenova.
\emph{Transformer-Based Language Models for Bulgarian}.
Proceedings of RANLP 2023 (Recent Advances in Natural Language Processing), 2023.

\bibitem{paev2024bulgarianllm}
N. Paev, K. Simov, P. Osenova, S. Petrova.
\emph{Large Language Models for Bulgarian}.
CLaDA-BG Conference, 2024.

\bibitem{omotoso2025bgem3}
A. Omotoso, H. Shopeju, A. M. Joshua, S. Oni.
\emph{Improving BGE-M3 Multilingual Dense Embeddings for Nigerian Low Resource Languages}.
Proceedings of the 9th Widening NLP Workshop (WiNLP), pp. 224--229, 2025.

\bibitem{nzeyimana2025kinyacolbert}
A. Nzeyimana, A. N. Rubungo.
\emph{KinyaColBERT: A Lexically Grounded Retrieval Model for Low-Resource Retrieval-Augmented Generation}.
arXiv:2507.03241, 2025.

\bibitem{navasardyan2026embeddings}
Z. Navasardyan, S. Bughdaryan, B. Minasyan, H. Davtyan.
\emph{Less is More: Adapting Text Embeddings for Low-Resource Languages with Small Scale Noisy Synthetic Data}.
arXiv:2603.22290, 2026.

\bibitem{alexandrov2024bggpt}
A. Alexandrov, V. Raychev, D. I. Dimitrov, C. Zhang, M. Vechev, K. Toutanova.
\emph{BgGPT 1.0: Extending English-centric LLMs to other languages}.
arXiv:2412.10893, 2024.

\bibitem{chen2024bgem3}
J. Chen, S. Xiao, P. Zhang, K. Luo, D. Lian, Z. Liu.
\emph{M3-Embedding: Multi-Linguality, Multi-Functionality, Multi-Granularity Text Embeddings Through Self-Knowledge Distillation}.
Findings of the Association for Computational Linguistics: ACL 2024, 2024. arXiv:2402.03216.

\bibitem{wang2024multilinguale5}
L. Wang, N. Yang, X. Huang, L. Yang, R. Majumder, F. Wei.
\emph{Multilingual E5 Text Embeddings: A Technical Report}.
arXiv:2402.05672, 2024.

\bibitem{reimers2020multilingual}
N. Reimers, I. Gurevych.
\emph{Making Monolingual Sentence Embeddings Multilingual using Knowledge Distillation}.
Proceedings of the 2020 Conference on Empirical Methods in Natural Language Processing (EMNLP), pp. 4512--4525, 2020.

\bibitem{robertson2009bm25}
S. Robertson, H. Zaragoza.
\emph{The Probabilistic Relevance Framework: BM25 and Beyond}.
Foundations and Trends in Information Retrieval, 3(4):333--389, 2009.

\bibitem{cormack2009rrf}
G. V. Cormack, C. L. A. Clarke, S. Büttcher.
\emph{Reciprocal Rank Fusion Outperforms Condorcet and Individual Rank Learning Methods}.
Proceedings of the 32nd International ACM SIGIR Conference on Research and Development in Information Retrieval, pp. 758--759, 2009.

\bibitem{nogueira2019passage}
R. Nogueira, K. Cho.
\emph{Passage Re-ranking with BERT}.
arXiv:1901.04085, 2019.

\bibitem{liu2024lostmiddle}
N. F. Liu, K. Lin, J. Hewitt, A. Paranjape, M. Bevilacqua, F. Petroni, P. Liang.
\emph{Lost in the Middle: How Language Models Use Long Contexts}.
Transactions of the Association for Computational Linguistics, 12:157--173, 2024.

\bibitem{zhang2025qwen3embedding}
Y. Zhang и др.
\emph{Qwen3 Embedding: Advancing Text Embedding and Reranking Through Foundation Models}.
arXiv:2506.05176, 2025.

\bibitem{vera2025embeddinggemma}
H. Schechter Vera и др.
\emph{EmbeddingGemma: Powerful and Lightweight Text Representations}.
arXiv:2509.20354, 2025.

\bibitem{holtzman2020degeneration}
A. Holtzman, J. Buys, L. Du, M. Forbes, Y. Choi.
\emph{The Curious Case of Neural Text Degeneration}.
International Conference on Learning Representations (ICLR), 2020. arXiv:1904.09751.

\end{thebibliography}

\section*{ПРИЛОЖЕНИЯ}
% TODO: авторски код на Python по компоненти (векторно представяне и индексиране,
% извличане, генериране, демонстрационно приложение), добавян поетапно с
% напредването на реализацията -- по аналогия с номерираните приложения в
% примерната дипломна работа.

\end{document}
